% !TEX program = pdflatex
\documentclass[11pt,a4paper]{article}

\usepackage[margin=1in]{geometry}
\usepackage{amsmath,amssymb,amsthm,mathtools}
\usepackage[T1]{fontenc}
\usepackage{lmodern}
\usepackage{microtype}
\usepackage{booktabs}
\usepackage{array}
\usepackage{enumitem}
\usepackage{hyperref}
\usepackage{xcolor}

\hypersetup{
    colorlinks=true,
    linkcolor=blue,
    urlcolor=blue,
    citecolor=blue
}

\newtheorem{theorem}{Theorem}[section]
\newtheorem{proposition}[theorem]{Proposition}
\newtheorem{lemma}[theorem]{Lemma}
\newtheorem{corollary}[theorem]{Corollary}
\theoremstyle{definition}
\newtheorem{definition}[theorem]{Definition}
\theoremstyle{remark}
\newtheorem{remark}[theorem]{Remark}

\newcommand{\Miss}{\mathcal M}
\newcommand{\Hit}{\mathcal H}
\newcommand{\Mp}[1]{M_{#1}}
\newcommand{\Hp}[1]{H_{#1}}
\newcommand{\bbN}{\mathbb N}
\newcommand{\bbZ}{\mathbb Z}
\newcommand{\ip}{\le_p}
\newcommand{\pdigits}[2]{#1_0,#1_1,\ldots,#1_{#2}}

\title{From Lucas Digit Order to Local Digit Calculus:\\
Exact MISS Prefixes, Carry Gaps, and Path-Independent Transitions}
\author{IamLupo}
\date{14 September 2026}

\begin{document}
\maketitle

\begin{abstract}
This paper records the developments obtained after the previous arithmetic-sequence/Lucas-theorem paper.  The central object remains the digitwise MISS predicate arising from Lucas' theorem, but the investigation has now been reorganized around the exact combinatorics of the set
\[
\Miss_m=\{x\ge 0:x\ip m\}.
\]
Writing the base-$p$ digits of $m$ as $m=\sum_i m_i p^i$, the MISS set is shown to be a finite mixed-radix system.  This yields an exact rank parametrization, an explicit carry-gap classification, exact prefix-count identities, and a closed formula for the MISS prefix count.  The paper then develops the effect of an elementary no-carry digit update $m\mapsto m+p^r$.  The newly created MISS values form one explicit digit slice, producing an exact local transition formula.  Independent computations verify the prefix formula on 1,355,005 cases and the large arbitrary-state local transition on 308,952 valid cases, including safe 64-bit states, transitions crossing $2^{64}$, and updates that create a new highest digit.

The resulting structure can be viewed as a discrete calculus on base-$p$ digit space: elementary digit increments have explicit finite differences, and experiments show path independence under commuting no-carry updates.  The computational machinery is therefore substantially more mature than in the previous paper.  The remaining factoring problem is identified more sharply: the established framework efficiently measures the digit structure, but a factor-sensitive invariant linking this digit calculus to an unknown factor of a composite $N=pq$ has not yet been proved.  Several concrete routes toward such an invariant are formulated.
\end{abstract}

\section{Scope and relation to the previous work}

The preceding paper developed an arithmetic-sequence construction whose binomial coefficients can be filtered using Lucas' theorem.  The computational program subsequently moved away from repeated enumeration and toward an intrinsic description of the resulting HIT/MISS set.  The key observation is that, once a prime base $p$ is fixed and a reference integer $m$ is given, the relevant predicate is simply
\[
 x\ip m,
\]
meaning that every base-$p$ digit of $x$ is bounded by the corresponding digit of $m$.

The present paper therefore treats the digitwise order itself as the primary object.  The developments recorded here are:
\begin{enumerate}[label=(\arabic*)]
    \item a mixed-radix parametrization of all MISS values;
    \item an exact rank/successor description;
    \item a complete carry-gap classification;
    \item exact HIT/MISS prefix formulae;
    \item a corrected closed MISS-prefix formula;
    \item an elementary no-carry digit-transition theorem;
    \item path independence for commuting digit updates;
    \item independent computational validation designed specifically to separate mathematical errors from implementation-width errors;
    \item a clearer statement of what remains to be shown before the framework can become a factoring method.
\end{enumerate}

The numerical experiments below are evidence for the stated formulae and implementations; they are not substitutes for proofs.

\section{Digitwise order and the MISS set}

Let $p$ be prime and write
\[
 m=\sum_{i=0}^{d} m_i p^i,
 \qquad 0\le m_i<p.
\]
For another nonnegative integer
\[
 x=\sum_{i\ge0}x_i p^i,
\]
we write
\[
 x\ip m
 \quad\Longleftrightarrow\quad
 x_i\le m_i\text{ for every }i.
\]
The MISS and HIT sets relative to $m$ are
\[
\Miss_m=\{x\ge0:x\ip m\},
\qquad
\Hit_m=\{x\ge0:x\not\ip m\}.
\]
In the Lucas-theoretic applications, these are exactly the zero/nonzero binomial-coefficient regimes for the corresponding repeated-coefficient construction.

\begin{proposition}[Product formula for the MISS set]
The number of MISS values in the interval $[0,m]$ is
\[
 |\Miss_m|=
 \prod_{i=0}^{d}(m_i+1).
\]
\end{proposition}
\begin{proof}
A MISS value is obtained by independently choosing each digit $x_i$ from the finite set $\{0,1,\ldots,m_i\}$.  The choices are independent across digit positions, giving the product.
\end{proof}

The key point is that the set is not merely finite and multiplicative in size.  It has an explicit ordered coordinate system.

\section{Mixed-radix rank and successor structure}

Define the mixed-radix weights
\[
 W_0=1,
 \qquad
 W_i=\prod_{j<i}(m_j+1)
 \quad(i\ge1).
\]
For $x\in\Miss_m$, define
\[
 R_m(x)=\sum_i x_iW_i.
\]

\begin{theorem}[Exact rank parametrization]
The map
\[
 R_m:\Miss_m\longrightarrow
\left\{0,1,\ldots,
\prod_i(m_i+1)-1\right\}
\]
is a bijection.
\end{theorem}
\begin{proof}
The digits $x_i$ are precisely the mixed-radix digits with radices $m_i+1$.  Standard mixed-radix uniqueness gives both injectivity and surjectivity.
\end{proof}

Let
\[
 0=M_0<M_1<\cdots<M_R=m
\]
be the MISS values in ordinary integer order.  The rank map identifies the ordered MISS set with a consecutive integer interval in rank space.  Consequently, the successor operation has the usual mixed-radix form: locate the first digit below its maximum, increment it by one, and reset all lower digits to zero.

This provides an exact description of the intervals containing HIT values:
\[
 [M_k+1,M_{k+1}-1].
\]
Thus the entire HIT set is the complement of the mixed-radix successor skeleton.

\section{Carry gaps}

Suppose the successor from $M_k$ to $M_{k+1}$ carries at digit $r$.  The lower digits $0,1,\ldots,r-1$ reset, while digit $r$ increments.  The resulting gap length is
\[
 G_r
 =p^r-1-\sum_{i<r}m_i p^i.
\]
The number of gaps of this type is
\[
 C_r
 =m_r\prod_{i>r}(m_i+1).
\]

\begin{proposition}[Gap count]
The total number of MISS-to-MISS gaps is
\[
\sum_r C_r
 =
 \prod_i(m_i+1)-1.
\]
\end{proposition}
\begin{proof}
Every nonterminal MISS state has exactly one successor.  In mixed-radix coordinates, a carry at $r$ occurs precisely when all lower digits are maximal and digit $r$ is not maximal.  Counting the choices above $r$ and the possible current values of digit $r$ gives $C_r$.
\end{proof}

Summing the gap lengths gives the total number of HIT values up to $m$:
\[
 \sum_r C_rG_r
 =m+1-\prod_i(m_i+1).
\]
Hence
\[
 |\Hit_m\cap[0,m]|\;=
 m+1-|\Miss_m|.
\]

\section{Prefix counting}

Define the MISS prefix count
\[
 M_m(y)=
 \#\{x:0\le x\le y,\ x\ip m\}
\]
and the HIT prefix count
\[
 H_m(n)=
 \#\{x:0\le x<n,\ x\not\ip m\}.
\]
Then, for $n>0$,
\[
 H_m(n)=n-M_m(n-1).
\]
The prefix formulation is more useful than direct enumeration because it permits logarithmic-in-$m$ evaluation in the number of base-$p$ digits.

\subsection{Closed MISS-prefix formula}

Assume $y<m$, and let $h$ be the largest index for which
\[
 y_h\ne m_h.
\]
Because $y<m$ numerically, one has $y_h<m_h$.  Let
\[
 W_i=\prod_{j<i}(m_j+1).
\]
Then
\begin{theorem}[Closed MISS-prefix formula]
For $y<m$,
\[
 \boxed{
 M_m(y)=
 \sum_{i>h}y_iW_i
 +y_hW_h
 +M_{m_{<h}}(y_{<h})
 }
\]
where
\[
 m_{<h}=\sum_{i<h}m_ip^i,
 \qquad
 y_{<h}=\sum_{i<h}y_ip^i.
\]
For $y\ge m$,
\[
 M_m(y)=\prod_i(m_i+1).
\]
\end{theorem}
\begin{proof}
Fix the most significant digit $h$ where $y$ and $m$ differ.  All digits above $h$ must agree with $y$ in any admissible $x\le y$.  For every higher-prefix choice counted by $y_i$ at level $i>h$, the lower digits are unrestricted subject only to the digitwise bound from $m$, producing a block of size $W_i$.  At digit $h$, choosing $x_h<y_h$ gives $y_h$ complete lower blocks, each of size $W_h$.  The remaining case $x_h=y_h$ is exactly the lower prefix problem $x_{<h}\ip m_{<h}$ and $x_{<h}\le y_{<h}$.  Summing these disjoint contributions gives the formula.
\end{proof}

The case in which $y$ has implicit leading zero digits is handled automatically by padding missing digits with zero.  This is important for updates that create a new highest digit.

\section{Elementary no-carry digit updates}

Consider an elementary update at digit $r$:
\[
 m'=m+p^r,
 \qquad m_r<p-1.
\]
This is a no-carry update.  The digits below and above $r$ are unchanged and
\[
 m'_r=m_r+1.
\]

\begin{theorem}[Exact new MISS slice]
For a no-carry update $m'=m+p^r$, the newly created MISS values are precisely
\[
 \mathcal S_r(m)=
 \left\{
 x:\ x_r=m_r+1,
 \quad x_i\le m_i\text{ for }i\ne r
 \right\}.
\]
Consequently,
\[
 M_{m'}(y)-M_m(y)
 =\#\{x\le y:x\in\mathcal S_r(m)\}.
\]
\end{theorem}
\begin{proof}
The only changed digit bound is at position $r$.  A value that was previously forbidden but becomes admissible must therefore violate the old bound only at digit $r$, and the new maximum there is exactly $m_r+1$.  All other digits remain bounded by their original values.
\end{proof}

This is the fundamental local transition identity.  It converts a global difference between two MISS-prefix functions into an explicit digit slice.

\subsection{Closed local-slice decomposition}

Write
\[
 y
 =H_y p^{r+1}+y_rp^r+L_y,
\qquad
 m
 =H_m p^{r+1}+m_rp^r+L_m.
\]
Let
\[
 B_r(m)=\prod_{i<r}(m_i+1).
\]
The newly created slice fixes digit $r$ to $m_r+1$, leaves the lower digits in a MISS block, and leaves the higher digits in a MISS prefix.  Therefore its prefix count splits into three regimes:
\begin{enumerate}[label=(\alph*)]
    \item higher prefix strictly below $H_y$;
    \item equal higher prefix, with $y_r<m_r+1$;
    \item equal higher prefix, with $y_r=m_r+1$ and a partial lower block.
\end{enumerate}
In the resulting closed formula, the lower complete block size is $B_r(m)$, the higher complete-block count is a MISS prefix in $H_m$, and the boundary term is another MISS prefix in $L_m$.

The important conceptual point is that an elementary digit change affects only one digit slice; no search through the entire interval $[0,m]$ is required.

\section{Path independence and discrete calculus}

Let $\Delta_r$ denote the no-carry forward difference
\[
 \Delta_rF(m)=F(m+p^r)-F(m).
\]
For distinct digits whose updates are simultaneously legal, the experiments support the commutation law
\[
 \Delta_r\Delta_sF
 =
 \Delta_s\Delta_rF.
\]
Repeated experiments strengthened this from two updates to arbitrary tested collections of updates and all tested permutations of the update order.

Equivalently, starting from a base state $m$ and applying a collection of independent elementary digit increments, the total change depends only on the final digit vector, not on the order of the path.

This motivates viewing the construction as a discrete calculus on the lattice of digit states.  The elementary local slices are discrete first derivatives, and the observed commutativity is the discrete analogue of path independence for a conservative field.

\begin{remark}
The computational evidence for path independence is stronger than needed to evaluate a single prefix, but the proof-theoretic status is different: the local slice theorem already identifies the elementary increments as set differences, so path independence is fundamentally explained by set-theoretic telescoping whenever the updates remain within a common digitwise partial order.  The experiments serve as validation of the explicit formulas used to calculate those increments.
\end{remark}

\section{Computational validation}

A major part of this stage of the project was separating mathematical statements from implementation errors.  Several intermediate scripts produced false failures caused by integer-width truncation or incorrect reference-prefix implementations.  The final validation protocol was redesigned to use an independent digit-DP reference and to test the dangerous regimes separately.

\subsection{Prefix validation}

Experiment 288 compared the corrected closed MISS-prefix formula against an independent most-significant-digit dynamic program.  The results were:
\[
\begin{array}{lrr}
\toprule
\text{test family} & \text{cases} & \text{failures}\\
\midrule
\text{small exhaustive} & 1,255,005 & 0\\
\text{large random} & 100,000 & 0\\
\midrule
\text{total} & 1,355,005 & 0\\
\bottomrule
\end{array}
\]
This validates the prefix formula over both small exhaustive states and large arbitrary states.

\subsection{Independent local-slice validation}

Experiment 284 used a different strategy: it defined the newly created MISS slice directly from its digit conditions and compared it to both the global MISS-prefix difference and the proposed closed local-slice formula.  The result was
\[
\begin{array}{lrr}
\toprule
&\text{cases}&\text{failures}\\
\midrule
\text{deterministic}&201,828&0\\
\text{boundary}&2,135&0\\
\text{random}&4,406&0\\
\midrule
\text{total}&208,369&0\\
\bottomrule
\end{array}
\]
with all three comparisons passing:
\[
\text{global delta}
=
\text{explicit new slice}
=
\text{closed local slice}.
\]

\subsection{Large arbitrary-current validation}

Experiment 289 then retested the local transition at genuinely large arbitrary states using the corrected prefix formula against the independent digit-DP reference.  The test was deliberately partitioned into prime ranges and arithmetic-width regimes:
\[
\begin{array}{lrrr}
\toprule
\text{group} & \text{cases} & \text{failures} & \text{new-digit cases}\\
\midrule
\text{edge cases}&8,952&0&2,912\\
\text{low primes}&100,000&0&1\\
\text{middle primes}&100,000&0&5,553\\
\text{high primes}&100,000&0&18,732\\
\midrule
\text{total}&308,952&0&27,198\\
\bottomrule
\end{array}
\]
The high-prime group contained 12,158 transitions crossing the unsigned 64-bit boundary, and all of them passed.  Across all groups,
\[
\text{safe64 failures}=0,
\qquad
\text{overflow64 failures}=0,
\qquad
\text{new-digit failures}=0.
\]
Thus the local formula was tested not only in conventional 64-bit-safe states but also in precisely the cases that had previously exposed implementation-width mistakes.

\section{What has actually been established}

The present stage can be summarized by the following chain:
\[
\boxed{
\begin{array}{c}
\text{Lucas/digitwise order}\\
\downarrow\\
\text{mixed-radix MISS set}\\
\downarrow\\
\text{exact rank and successor}\\
\downarrow\\
\text{carry-gap decomposition}\\
\downarrow\\
\text{exact MISS/HIT prefixes}\\
\downarrow\\
\text{closed MISS-prefix formula}\\
\downarrow\\
\text{elementary digit slices}\\
\downarrow\\
\text{path-independent digit transitions}
\end{array}}
\]
This is a substantially stronger structural description than the original term-by-term formulation.  The computational object can now be manipulated through local digit coordinates rather than through an exponentially or linearly large collection of individual integers.

There is also an important algorithmic consequence.  Direct enumeration of all $x<n$ is unnecessary for prefix evaluation: the closed formula operates on the base-$p$ digit expansion of $m$ and $n$, so the cost grows essentially with the number of digits rather than with the numerical magnitude of the interval.

\section{What has \emph{not} yet been established}

The present work does \emph{not} yet constitute a factoring algorithm for a general composite $N$.  The missing ingredient is a theorem connecting the digit calculus to an unknown divisor of $N$.

For a semiprime
\[
 N=pq,
 \qquad p\le q,
\]
the current framework can calculate many quantities attached to a chosen base-$p$ digit system, but a successful factoring mechanism would require a factor-sensitive observable.  Typical desired forms include
\[
 F(N,x)\equiv0\pmod p,
 \qquad
 F(N,x)\not\equiv0\pmod q,
\]
so that
\[
 \gcd(F(N,x),N)=p,
\]
or a threshold function whose transition location is provably related to $p$.

The distinction is important.  The present paper establishes an exact and efficiently evaluable \emph{measurement framework}; it does not yet prove that this measurement leaks a nontrivial factor for arbitrary $N$.

\section{Potential routes toward factor extraction}

The local digit calculus suggests several concrete research directions.

\subsection{Factor-sensitive finite differences}

The first possibility is to construct an $N$-dependent function $F$ from elementary differences such as
\[
 \Delta_rF,
 \qquad
 \Delta_r\Delta_sF,
\]
and search for expressions whose residues behave asymmetrically modulo the two prime factors of $N$.

The path-independent structure is particularly attractive here because it permits algebraic combinations of elementary transitions without concern about the order used to evaluate them.

\subsection{Threshold localization}

The earlier arithmetic-sequence work suggested a chain of the form
\[
\text{arithmetic construction}
\rightarrow
\binom{m}{d}
\rightarrow
\text{Lucas criterion}
\rightarrow
\text{threshold predicate}
\rightarrow
\text{binary localization}.
\]
The current prefix machinery makes the localization stage computationally cheap.  The remaining question is the deeper one: can the threshold be proved to occur at a quantity determined by an unknown factor of $N$?

\subsection{GCD extraction from polynomial residues}

A related route is to construct a polynomial or integer-valued function $F_N(x)$ for which one can prove that a strategically selected argument $x$ produces
\[
 F_N(x)\equiv0\pmod p
\]
while avoiding the same congruence modulo $q$.  The experimentally observed usefulness of gcd tests in the earlier polynomial experiments makes this a particularly concrete target.

\subsection{Second and higher discrete derivatives}

Because elementary updates commute, the quantities
\[
 \Delta_r\Delta_sF,
 \quad
 \Delta_r\Delta_s\Delta_tF,
 \quad\ldots
\]
are natural invariants of the digit-state geometry.  These higher differences may cancel nuisance terms in a local slice formula and leave a simpler factor-sensitive quantity.  This has not yet been developed into a proof of factor recovery, but it is a promising algebraic direction.

\section{Discussion}

The main conceptual change since the previous paper is that the original arithmetic construction is no longer being treated primarily as a sequence of isolated binomial evaluations.  Instead, it has been converted into a finite combinatorial geometry indexed by base-$p$ digits.

The MISS set is mixed-radix.  Its complement consists of explicit gaps.  Prefix counts admit closed formulae.  Elementary digit changes create localized slices.  Those slices commute under independent updates.  Consequently, the entire system admits a discrete potential/gradient interpretation.

This does not by itself solve factorization.  Nevertheless, it changes the nature of the remaining problem.  The question is no longer how to efficiently enumerate the underlying combinatorial objects; that part has an exact solution.  The question is now whether one can choose the digit-state, the base, or a combination of elementary derivatives so that the resulting observable is sensitive to one factor of a composite and not the other.

In this sense, the present work should be viewed as the completion of the \emph{combinatorial engine} needed for a possible factoring construction.  The factor-sensitive invariant remains the central open step.

\section{Conclusion}

The post-previous-paper developments establish a coherent algebraic chain from Lucas digitwise order to an exact local digit calculus.  The core identities are:
\[
|\Miss_m|=\prod_i(m_i+1),
\]
\[
\sum_r C_rG_r
=m+1-\prod_i(m_i+1),
\]
\[
H_m(n)=n-M_m(n-1)\quad(n>0),
\]
and, for the corrected explicit prefix formula,
\[
M_m(y)=
\sum_{i>h}y_iW_i
+y_hW_h
+M_{m_{<h}}(y_{<h}),
\]
followed by the local no-carry transition
\[
M_{m+p^r}(y)-M_m(y)
=\#\{x\le y:x_r=m_r+1,\ x_i\le m_i\ (i\ne r)\}.
\]

The computational program now supports these identities over more than a million independent prefix cases and hundreds of thousands of large arbitrary-state local updates, including deliberately difficult integer-width and new-digit regimes.

The next mathematical objective is therefore clear: construct and prove a factor-sensitive observable derived from this digit calculus.  If such an observable can be made asymmetric modulo the two prime factors of a semiprime $N=pq$, the final extraction step would be a standard gcd computation.  Until that link is proved, the present framework should be regarded as a strong structural and algorithmic foundation rather than a complete factoring algorithm.

\end{document}
